\chapter{2022:Maryna Viazovska}\label{ch2022marynaviazovska:maryna-viazovska-fields-medal-2022-a-comprehensive-report-on-her-work-and-impact}

\begin{center}\rule{0.5\linewidth}{0.5pt}\end{center}

\textbf{Prepared:} July 2026

\textbf{Subject:} Maryna Viazovska (b.\ December 2, 1984, Kyiv, Ukraine)

\textbf{Position:} Chair of Number Theory, \'{E}cole Polytechnique F\'ed\'erale de Lausanne (EPFL)

\textbf{Doctoral Advisor:} Don Zagier (University of Bonn), then Ph.\ Michael M\"uller (Bonn)

\textbf{Fields Medal Citation:} \emph{``For her work on the densest packing of equal spheres in eight dimensions, and for contributions to related extremal problems in several dimensions.''}

\begin{center}\rule{0.5\linewidth}{0.5pt}\end{center}

\begin{prereqs}
Number theory (modular forms, theta functions, Fourier analysis), geometry of numbers (lattices, sphere packings), harmonic analysis (Bessel functions, Hankel transforms).
\end{prereqs}

\begin{objectives}
Explain the sphere packing problem and its history; describe the E8 lattice and Leech lattice; understand the connection between modular forms and packing densities via the Cohn--Elkies criterion.
\end{objectives}

\section{Problem History}

The \textbf{sphere packing problem} asks: what is the maximum density of a packing of equal non-overlapping spheres in $\mathbb{R}^n$? In dimension 2, the hexagonal lattice achieves the optimal density $\pi/\sqrt{12} \approx 0.9069$ (Gauss, 1831). In dimension 3, the face-centered cubic (FCC) lattice is conjectured optimal (density $\pi/\sqrt{18} \approx 0.7405$), but this remains unproved (the Kepler conjecture, proved by Hales, 2005).

In dimension 8, the \textbf{E8 lattice} is a highly symmetric lattice constructed from the root system of the Lie group $E_8$. It achieves the density:

\[
\delta_8 = \frac{\pi^4}{384} \approx 0.2537
\]

In dimension 24, the \textbf{Leech lattice} achieves density:

\[
\delta_{24} = \frac{\pi^{12}}{12!} \approx 0.00193
\]

Conway and Sloane conjectured in the 1980s that the E8 and Leech lattices are optimal in their respective dimensions, based on computational evidence and the Levenshtein--Cohn--Mills bounds.

The \textbf{Cohn--Kumar bound} (2003) showed that if a lattice packing in $\mathbb{R}^n$ has density equal to $\delta_n$ (the Coxeter--Todd bound), then the lattice must be the E8 or Leech lattice. But this was merely a numerical check---it did not prove the bound itself.

The classical approach to sphere packing upper bounds uses the \textbf{Cohn--Elkies criterion} (2003): a function $\phi: \mathbb{R}^n \to \mathbb{R}$ satisfying $\widehat{\phi}(x) \le 0$ for $|x| \ge 1$, $\phi(x) \ge 0$ for all $x$, and $\widehat{\phi}(0) = \phi(0)$ provides an upper bound on the packing density. Finding such a function was the key challenge.

\section{How Was It Solved}

Viazovska's approach was to construct explicit functions satisfying the Cohn--Elkies criterion using modular forms.

\textbf{The modular form approach.} For a lattice $\Lambda \subset \mathbb{R}^n$, the theta function is:

\[
\Theta_\Lambda(z) = \sum_{x \in \Lambda} e^{2\pi i |x|^2 z}
\]

This is a modular form of weight $n/2$ for the full modular group $\operatorname{SL}_2(\mathbb{Z})$. The density of the lattice packing is determined by the first non-zero Fourier coefficient of $\Theta_\Lambda$.

\textbf{The construction of the magic function (2016).} Viazovska defined a modular form $f$ of weight 4 for $\operatorname{SL}_2(\mathbb{Z})$ such that $f(z) = 1 + a_1 q + a_2 q^2 + \cdots$ with $q = e^{2\pi i z}$. The key property is that $f(z)$ is \textbf{non-negative} on the upper half-plane and $f(i) = 0$. Using the theory of modular forms (dimension of spaces, valence formula), she constructed explicit coefficients.

The space of modular forms of weight 4 for $\operatorname{SL}_2(\mathbb{Z})$ is one-dimensional, spanned by the Eisenstein series $E_4(z) = 1 + 240q + 2160q^2 + \cdots$. However, Viazovska's function $f$ is not a modular form in the classical sense---it is a linear combination of Eisenstein series and cusp forms that satisfies weaker transformation properties. The key innovation is to work with a broader class of functions that include the Eisenstein series and have the required positivity properties.

The function $\phi(x) = \sum_{n \in \mathbb{Z}} f(n^2 x) |n|$ is then a radial function in $L^1(\mathbb{R}^8)$ such that:

\[
\widehat{\phi}(x) \le 0 \quad \text{for } |x| \ge 1, \qquad \widehat{\phi}(0) = \frac{\pi^4}{384}
\]

By the Cohn--Elkies criterion, any such function provides an upper bound on the packing density. Since $\widehat{\phi}(0) = \delta_8$, this proves the E8 lattice is optimal.

\textbf{The 24-dimensional case (2017).} With Cohn, Kumar, Miller, Radchenko, and Viazovska, the same approach was applied to the Leech lattice using weight 12 modular forms, completing the proof for dimension 24.

\textbf{The Fourier transform of radial functions.} Viazovska and collaborators introduced a new method for computing the Fourier transform of radial functions in $L^1(\mathbb{R}^n)$, using the Hankel transform and integral identities for modified Bessel functions. This provided the technical foundation for both the 8-dimensional and 24-dimensional proofs.

The Fourier transform of a radial function $\phi(x) = f(|x|)$ in $\mathbb{R}^n$ can be expressed in terms of the Hankel transform:
\[
\widehat{\phi}(\xi) = \frac{2\pi^{n/2}}{|\xi|^{(n-2)/2}} \int_0^\infty f(r) J_{(n-2)/2}(2\pi r |\xi|) r^{n/2} \, dr
\]
where $J_\nu$ is the Bessel function of the first kind. Viazovska's key innovation was to use integral identities for modified Bessel functions $I_\nu$ to compute the Fourier transform of functions of the form $\phi(x) = \sum_{n \in \mathbb{Z}} f(n^2 |x|^2)$, where $f$ is a modular form. This allowed her to verify the Cohn--Elkies criterion by showing that $\widehat{\phi}$ satisfies the required sign conditions.

The technical difficulty lies in the fact that the Fourier transform of a sum over a lattice is itself a sum over the dual lattice (by the Poisson summation formula), and controlling the tail of this sum requires precise estimates on the Bessel functions. Viazovska's method provides these estimates using the theory of modular forms and the properties of the Bessel function near the origin and at infinity.

\section{Impact of This Discovery}

The sphere packing problem, open since Kepler in 1611, was solved in its two most important cases. Viazovska's method is completely elementary---it uses only modular forms, Fourier analysis, and the Cohn--Elkies criterion, with no heavy machinery.

The explicit construction of the ``magic function'' has applications to coding theory (the E8 and Leech lattices are optimal error-correcting codes), signal processing, and the theory of sphere packings.

The Fourier transform techniques developed by Viazovska and collaborators have found applications in analytic number theory, in the theory of positive-definite functions, and in harmonic analysis.

The \textbf{Cohn--Elkies criterion} (2003) provides a necessary and sufficient condition for optimality of a sphere packing. A function $\phi: \mathbb{R}^n \to \mathbb{R}$ satisfying $\widehat{\phi}(x) \le 0$ for $|x| \ge 1$, $\phi(x) \ge 0$ for all $x$, and $\widehat{\phi}(0) = \phi(0)$ provides an upper bound on the packing density $\delta_n \le \phi(0)/2^n$. If equality holds for a lattice $\Lambda$ with minimal distance 1, then $\Lambda$ is optimal. The challenge is to construct such a function, which Viazovska accomplished for dimensions 8 and 24.

The \textbf{E8 lattice} is a highly symmetric lattice in $\mathbb{R}^8$ with 240 minimal vectors (kissing number 240). It can be constructed as the set of points in $\mathbb{R}^8$ whose coordinates are all integers or all half-integers and whose sum is even. The theta function of the E8 lattice is a modular form of weight 4 for $\operatorname{SL}_2(\mathbb{Z})$, and its Fourier expansion is
\[
\Theta_{E_8}(z) = 1 + 240q + 2160q^2 + \cdots, \quad q = e^{2\pi i z}.
\]
The packing density of the E8 lattice is $\delta_8 = \pi^4/384 \approx 0.2537$, and Viazovska's proof that this is optimal uses the modular form $f(z) = E_4(z) - 3\Delta(z)/\Delta(z)$, where $E_4$ is the Eisenstein series of weight 4 and $\Delta$ is the modular discriminant.

The \textbf{Leech lattice} is a highly symmetric lattice in $\mathbb{R}^{24}$ with 196560 minimal vectors (kissing number 196560). It can be constructed using the Golay code, an error-correcting code of length 24. The theta function of the Leech lattice is a modular form of weight 12 for $\operatorname{SL}_2(\mathbb{Z})$, and its packing density is $\delta_{24} = \pi^{12}/12! \approx 0.00193$. The proof that the Leech lattice is optimal uses a weight 12 modular form, following the same strategy as the 8-dimensional case but with more intricate technical details.

\section{Modern Views and Improvements}

Viazovska continues to develop her approach to sphere packing, with current work on packing in dimension 4 (where the D4 lattice is conjectured optimal) and on bounds in general dimensions.

Her recent work with Radchenko on the theory of \textbf{period functions} of Maass forms connects to the classical theory of the Riemann zeta function and has applications to the functional equation. The period functions are Fourier coefficients of Maass forms, and they satisfy a functional equation that is closely related to the functional equation of the Riemann zeta function. Viazovska and Radchenko have developed new methods for computing period functions that have applications to analytic number theory and the theory of automorphic forms.

The method has been extended to other extremal problems: the kissing number problem (how many unit spheres can simultaneously touch a central unit sphere) was resolved in dimensions 8 and 24 using related techniques. The kissing number of the E8 lattice is 240, and the kissing number of the Leech lattice is 196560. These results follow from the Cohn--Elkies criterion applied to the kissing number problem, and the proof uses modular forms in a similar way to the sphere packing problem.

The connection between sphere packing and coding theory is another active direction. The E8 and Leech lattices are optimal error-correcting codes, and the sphere packing density gives a lower bound on the rate of the code. Viazovska's work has implications for coding theory by providing the optimal sphere packing densities in dimensions 8 and 24.

The development of new techniques for computing Fourier transforms of radial functions has applications beyond the sphere packing problem. The integral identities for modified Bessel functions that Viazovska and collaborators introduced have been applied to the study of positive-definite functions, to the theory of positive polynomials, and to the analysis of theta functions. The connection between modular forms and Fourier analysis, as demonstrated in Viazovska's work, continues to produce new results in number theory and harmonic analysis.

\begin{center}\rule{0.5\linewidth}{0.5pt}\end{center}

\section*{Summary}
\addcontentsline{toc}{section}{Summary}

Viazovska proved that the E8 lattice is the densest sphere packing in dimension 8 using a novel construction of modular forms and the Cohn--Elkies criterion. The same methods, extended with collaborators, proved the Leech lattice is optimal in dimension 24.
